\appendix
\section{Repositorio y reproducibilidad}
\label{sec:repositorio}

Todo el material está en \url{\repositorio}: el Verilog (\sig{hdl/}), los bancos (\sig{sim/}), la terminal
(\sig{pc/}), los proyectos de Quartus con sus reportes y archivos \sig{.sof} (\sig{quartus\_tri\_*/}), el proyecto de
Active-HDL, el layout con su GDS, la simulación de misión y las figuras (\sig{modelo/}) y las notas con cada dato y su
fuente (\sig{investigacion/}). \sig{GUIA\_LABORATORIO.md} describe el cableado y el guion de la demostración.

\begin{table}[!htb]
\centering\footnotesize
\caption{Cómo se reproduce cada resultado.}
\label{tab:reproducir}
\rowcolors{2}{SoftGray}{white}
\begin{tabularx}{\textwidth}{@{}>{\raggedright\arraybackslash}p{6.6cm} >{\ttfamily\raggedright\arraybackslash}X@{}}
\toprule
\rowcolor{TableHead} \textbf{Resultado} & \normalfont\textbf{Orden, desde la raíz del repositorio} \\
\midrule
Bancos de prueba y los nueve escenarios & pwsh -File sim\textbackslash simular.ps1 -Completa \\
Terminal contra las tarjetas simuladas & pwsh -File sim\textbackslash probar\_terminal.ps1 \\
Síntesis de las cuatro variantes & pwsh -File sintetizar\_red.ps1 \\
Carga de las tarjetas y demostración & CARGAR\_RED.cmd y DEMO.cmd \\
Layout en SKY130 & layout/correr\_flujo.sh \\
Simulación de misión y documentos & python modelo/experimentos.py; pwsh -File compilar.ps1 \\
\bottomrule
\end{tabularx}
\end{table}
